\documentclass[10pt,a4paper]{amsart}
\usepackage[margin=19mm]{geometry}
\usepackage{amsmath,amssymb,mathtools}
\usepackage[scr=rsfs,cal=euler]{mathalfa}
\usepackage{xfrac}
\usepackage[unicode,hidelinks,bookmarks=false]{hyperref}
\newcommand{\ie}{i.e.\ }
\newcommand{\cf}{cf.\ }
\newcommand{\resp}{respectively\ }
\newcommand{\Subsection}{\subsection}
\providecommand{\nobreakdash}{\nobreak\hbox{-}}
\setlength{\emergencystretch}{2em}
\multlinegap0pt
\newtheorem{theorem}{Theorem}
\newtheorem{lemma}[theorem]{Lemma}
\newcommand{\De}{\Delta}
\newcommand{\Om}{\Omega}
\newcommand{\ba}{\beta}
\newcommand{\pa}{\partial}
\newcommand{\sr}{\mathbb{R}}
\newcommand{\si}{\sigma}
\newcommand{\na}{\nabla}
\newcommand{\omc}{(\overline{\Om})}
\newcommand{\bdry}{\mathbb{R}^N \setminus \Omega}
\newcommand{\et}{\tilde{e}_0}
\title[Lemma 6.4]{Multiple positive solutions to a perturbed Gelfand problem involving mixed local-nonlocal operators and singular nonlinearity: the Hopf-type strong comparison principle (Lemma 6.4)}
\author{Sarbani Pramanik}
\date{Source: arXiv:2411.19694v5 (CC BY 4.0)}
\begin{document}
\maketitle

\noindent\emph{Source and licence.} Multiple positive solutions to a perturbed Gelfand problem involving mixed local-nonlocal operators and singular nonlinearity, by Sarbani Pramanik. Version arXiv:2411.19694v5
(CC BY 4.0), distributed by arXiv under the Creative Commons Attribution 4.0 International licence,
http://creativecommons.org/licenses/by/4.0/, as recorded in the arXiv licence metadata for this exact
version and shown on the abstract page https://arxiv.org/abs/2411.19694v5. Full licence notice:
\texttt{licenses/arxiv-2411.19694-CC-BY-4.0.txt}.\par\smallskip

\noindent\textit{Editorial scope.} Counted unit: the article's Lemma 6.4 (source0.tex lines 712--714, printed as Lemma 6.4): for $s\in[\frac12,1)$, two positive $C^{1,\alpha}(\overline{\Om})$ functions satisfying the singular mixed local-nonlocal system with $0\le f_1\le f_2$, $f_1\not\equiv 0$, $f_1\not\equiv f_2$ obey $\frac{\pa u_2}{\pa\nu}<\frac{\pa u_1}{\pa\nu}<0$ on $\pa\Om$. The article's complete original proof of it is delivered with it (source0.tex lines 716--765, one \texttt{proof} environment of 50 lines, adjacent to the statement, with no gap and no deferral). No mathematics is written, repaired or re-derived and no retained line is substituted or re-flowed; the only adaptation anywhere in the package is stated next. Retained uncounted as the source-local context the proof needs, each exactly the statement of the result it invokes and nothing else: the linear Strong Comparison Principle with its displayed system (lines 174--184, printed as Theorem 1.3, source label \texttt{Hopf lemma\_linear case}); the comparison lemma that the auxiliary function $w$ satisfies (lines 611--613, printed as Lemma 6.1) together with the displayed differential inequality it produces (lines 621--623, source label \texttt{w\_equn}); the $C^{1,\alpha}$-extension lemma with its equation for $e_0$ (lines 650--659, printed as Lemma 6.2); the statement of the companion lemma for $s\in(0,\frac12)$ (lines 682--684, printed as Lemma 6.3, source label \texttt{hopf\_first}); and the article's complete proof of that companion lemma (lines 686--710), which is what defines the auxiliary functions $w$ and $v$ that the counted proof invokes, so that the counted argument can be followed inside this document without a gap. Every cross-reference of every retained line resolves inside this document.

One declared adaptation, and the only place where the delivered text differs from the article's own lines: the two attribution commands in the retained proof of Lemma 6.3 are replaced, one for one and with nothing else changed, by DOI hyperlinks carrying the same bibliographic details transcribed from the article's own bibliography file \texttt{reference.bib} in the e-print -- \texttt{\textbackslash href\{https://doi.org/10.1016/j.jde.2025.01.030\}\{...\}} for the work of Antonini and Cozzi (a Hopf lemma for mixed operators) and \texttt{\textbackslash href\{https://doi.org/10.1002/cpa.20153\}\{...\}} for the work of Silvestre (regularity for a fractional power of the Laplacian). Both substitutions are recorded in the item's reference mapping and in the package delivery evidence. Because those two commands were replaced and no other citation occurs in any retained line, the wrapper carries no bibliography and no transcribed bibliography entry; the article's own bibliography file is not redistributed.

Editorial formatting only, in addition: an AMS \texttt{amsart} preamble that restates the article's own declarations and macros used by the retained lines -- the environments \texttt{theorem} and \texttt{lemma}, kept on one shared counter exactly as in the article but without the article's section-based prefix, and the macros \texttt{\textbackslash De}, \texttt{\textbackslash Om}, \texttt{\textbackslash ba}, \texttt{\textbackslash pa}, \texttt{\textbackslash sr}, \texttt{\textbackslash si}, \texttt{\textbackslash na}, \texttt{\textbackslash omc}, \texttt{\textbackslash bdry} and \texttt{\textbackslash et}, copied from the article's own preamble (source0.tex lines 49--69) -- and the article's own \texttt{amsart} class. The retained environments print in this document as Theorem 1, Lemma 2, Lemma 3, Lemma 4 and Lemma 5 in order of appearance, whereas the article prints the same items as Theorem 1.3, Lemma 6.1, Lemma 6.2, Lemma 6.3 and Lemma 6.4; the fourteen numbered displays of the retained lines print here as (1) to (14) in order of appearance, whereas the article numbers its equations per section.

The complete native source of the article as distributed in the arXiv e-print for this exact version is bundled unchanged as \texttt{sources/169340/source0.tex}: it is byte identical to the e-print file \texttt{three\_sol\_arxiv.tex} except that the pipeline's mechanical concatenation prepends one header comment line naming it. No publisher class or style file is shipped in the e-print and none is redistributed or used.
\begin{theorem}[\textbf{Strong Comparison Principle}]\label{Hopf lemma_linear case}
    Let $u_1, u_2 \in C^{1,\alpha}(\overline{\Om})$, for some $\alpha\in(0,1)$, $u_1, u_2 > 0$ in $\Om$ satisfy 
    \begin{equation}\label{scp_equns}
        \begin{aligned}
            -\Delta u_1 + (-\De)^s u_1 - \frac{1}{u_1^{\ba}} = f_1 \text{ in } \Om\\
            -\Delta u_2 + (-\De)^s u_2 - \frac{1}{u_2^{\ba}} = f_2 \text{ in } \Om\\
            u_1 =u_2 =0 \text{ in } \bdry
        \end{aligned}
    \end{equation}
    where $f_1,f_2$ are continuous functions with $0 \leq f_1 \leq f_2$, $f_1 \not\equiv 0$ and $f_1 \not\equiv f_2$ in $\Om$. Then $u_1 < u_2$ in $\Om$ and $\frac{\pa u_2}{\pa\nu} < \frac{\pa u_1}{\pa\nu} < 0$ on $\pa\Om$.
\end{theorem}

\begin{lemma}\label{strict comparison}
    Let $u_1, u_2 \in C^{1,\alpha}(\overline{\Om})$, for some $\alpha\in(0,1)$, $u_1, u_2 > 0$ in $\Om$ and they satisfy \eqref{scp_equns} with $f_1, f_2$ as in Theorem \ref{Hopf lemma_linear case}. Then $u_1<u_2$ in $\Om$.
\end{lemma}

    \begin{equation}\label{w_equn}
        -\Delta w + (-\De)^s w + \frac{k}{d(x)^{\ba +1}} w \geq 0 \text{ in } \Om.
    \end{equation}

\begin{lemma}\label{extension of e}
    Let $e_0 \in C^{1,\alpha} \omc$, $\alpha \in (0,1)$, be the unique solution of
    \begin{equation}\label{e_0 equation}
        \begin{aligned}
            -\Delta e_0 &=1 \text{ in } \Om\\
            e_0 &= 0 \text{ on } \pa\Om.
        \end{aligned}
    \end{equation}
    Then there exists an extension $\tilde{e}_0$ of $e_0$ such that $\tilde{e}_0 \in C^{1,\alpha} (\sr^N)$, $\tilde{e}_0 = e_0$ in $\overline{\Om}$ and $\tilde{e}_0 \leq 0$ in $\sr^N \setminus \Om$.
\end{lemma}

\begin{lemma}\label{hopf_first}
    Let $s \in \left(0, \frac{1}{2}\right)$, $u_1, u_2 \in C^{1,\alpha}(\overline{\Om})$, for some $\alpha\in(0,1)$, $u_1, u_2 > 0$ in $\Om$ and satisfy \eqref{scp_equns} with $f_1, f_2$ as in Theorem \ref{Hopf lemma_linear case}. Then $\frac{\pa u_2}{\pa\nu} < \frac{\pa u_1}{\pa\nu} < 0$ on $\pa\Om$.
\end{lemma}

\begin{proof}
    From the weak comparison principle and the Hopf lemma for mixed operators (Theorem 1.2, \href{https://doi.org/10.1016/j.jde.2025.01.030}{C. A. Antonini and M. Cozzi, Global gradient regularity and a Hopf lemma for quasilinear operators of mixed local-nonlocal type, J. Differential Equations 425 (2025), 342-382}), it follows that $\frac{\pa u_2}{\pa\nu} \leq \frac{\pa u_1}{\pa\nu} < 0$ on $\pa\Om$. To refine this result and obtain the strict inequality $\frac{\pa u_2}{\pa\nu} < \frac{\pa u_1}{\pa\nu}$ on $\pa\Om$, we define $w:= u_2 - u_1$. By lemma \ref{strict comparison}, $w > 0$ in $\Om$ and satisfies \eqref{w_equn}.

    Define $v(x):= \et (x) + e_0 (x)^{\gamma}$, where $\et$ is the extension given by lemma \ref{extension of e}, $e_0$ (by abuse of notation) denotes the function solving \eqref{e_0 equation} with zero extension in $\sr^N \setminus \Om$ and the exponent $\gamma\in (1,2)$ is to be chosen later. Then
    \begin{equation}\label{direct computation}
        -\Delta v + \frac{k}{d(x)^{\ba +1}} v = 1 + \gamma e_0^{\gamma -1} -\gamma (\gamma -1) e_0^{\gamma -2} |\na e_0|^2 + \frac{k}{d(x)^{\ba +1}} \left(\et + e_0^{\gamma}\right) \text{ in } \Om.
    \end{equation}
    Since $\et (x) = e_0 (x) \sim d(x)$ in $\Om$ and $\inf_{\pa\Om} \left| \frac{\pa e_0}{\pa\nu} \right| >0$, choosing $\eta>0$ small enough and $\gamma \in (1, 2- \ba)$, we obtain
    \begin{equation}\label{est_k1}
        -\Delta v + \frac{k}{d(x)^{\ba +1}} v \leq - \frac{k_1}{d(x)^{2-\gamma}} \text{ in } \Om_{\eta},
    \end{equation}
    where $\Om_{\eta} = \left\{ x\in \Om : d(x) < \eta \right\}$ and $k_1 >0$ is a constant. Moreover, since $v \in C^{1,\alpha} (\sr^N)$ and $s\in \left(0, \frac{1}{2}\right)$, by Proposition 2.6 of \href{https://doi.org/10.1002/cpa.20153}{L. Silvestre, Regularity of the obstacle problem for a fractional power of the Laplace operator, Comm. Pure Appl. Math. 60 (2007), 67-112}, there exists a constant $k_2 >0$ such that $\left|(-\De)^s v\right| \leq k_2$ in $\overline{\Om}_{\eta}$. Therefore, choosing $\eta$ smaller if required, we get
    \begin{equation}\label{v_equn}
        -\Delta v + (-\De)^s v + \frac{k}{d(x)^{\ba +1}} v \leq 0 \text{ in } \Om_{\eta}.
    \end{equation}
    Since by lemma \ref{strict comparison}, $w >0$ in $\Om$, we choose $\epsilon>0$ small enough so that $\epsilon v \leq w$ in $\Om \setminus \Om_{\eta}$. Then, from \eqref{w_equn} and \eqref{v_equn} we obtain
    \begin{equation}
        \begin{aligned}
            -\Delta w + (-\De)^s w + \frac{k}{d(x)^{\ba +1}} w &\geq 0 \text{ in } \Om_{\eta}\\
            -\Delta (\epsilon v) + (-\De)^s (\epsilon v) + \frac{k}{d(x)^{\ba +1}} (\epsilon v) &\leq 0 \text{ in } \Om_{\eta}\\
            w \geq \epsilon v \text{ in } \sr^N \setminus \Om_{\eta}.
        \end{aligned}
    \end{equation}
    Thus, by weak comparison principle, we have $w \geq \epsilon v = \epsilon \left(\et + e_0^{\gamma}\right)$ in $\Om_{\eta}$. Since $\et = e_0$ in $\Om$, we thus have $\frac{\pa w}{\pa\nu} \leq \frac{\pa \left(\epsilon e_0\right)}{\pa\nu} < 0$ or, equivalently $\frac{\pa u_2}{\pa\nu} < \frac{\pa u_1}{\pa\nu}$ on $\pa\Om$. This completes the proof.
\end{proof}

\begin{lemma}\label{hopf_second}
    Let $s \in \left[\frac{1}{2}, 1\right)$, $u_1, u_2 \in C^{1,\alpha}(\overline{\Om})$, for some $\alpha\in(0,1)$, $u_1, u_2 > 0$ in $\Om$ and satisfy \eqref{scp_equns} with $f_1, f_2$ as in Theorem \ref{Hopf lemma_linear case}. Then $\frac{\pa u_2}{\pa\nu} < \frac{\pa u_1}{\pa\nu} < 0$ on $\pa\Om$.
\end{lemma}

\begin{proof}
    We define the functions $w$ and $v$ as in lemma \ref{hopf_first}. Recall that they satisfy, respectively, equations \eqref{w_equn} and \eqref{est_k1} in $\Om_{\eta}$, for $\eta>0$ small and $\gamma \in \left(1, 2- \ba\right)$.

    \noindent\textbf{Claim:} For sufficiently small $\eta>0$, there exists $k_2 >0$ such that for all $x\in \Om_{\eta}$,
    \[\left|(-\De)^s v (x)\right| \leq \frac{k_2}{d(x)^{2s-1}}, \text{ for } s\in \left[\frac{1}{2},1\right).\]

    To establish this, we first estimate the Hessian of $v$. For any $i,j= 1,2, \dots,$ $N$, we have
    \begin{equation*}
        \frac{\pa^2}{\pa x_i \pa x_j} \left(e_0^{\gamma}\right) = \gamma \left( \gamma -1 \right) e_0^{\gamma-2} \frac{\pa e_0}{\pa x_i} \frac{\pa e_0}{\pa x_j} + \gamma e_0^{\gamma-1} \frac{\pa^2 e_0}{\pa x_i \pa x_j} \text{ in } \Om_{\eta}.
    \end{equation*}
    Since $1< \gamma <2$ and $e_0$ is the unique solution of equation \eqref{e_0 equation}, it follows that
    \begin{equation}\label{Hessian_estimate}
        \left|D^2v (x)\right| \leq \frac{k_3}{d(x)^{2-\gamma}} \leq \frac{k_3}{d(x)} \text{ in } \Om_{\eta}, \text{ for } \eta>0 \text{ small enough}.
    \end{equation}
    Next, for $x\in \Om_{\eta}$, we decompose the nonlocal term as
    \begin{align}\label{est_nonlocal}
        (-\De)^s v (x) &= -\frac{1}{2} \int_{|z| \leq \frac{1}{2}d(x)} \hspace{-5pt} \frac{v(x+z) + v(x-z) - 2v(x)}{|z|^{N+2s}}\ dz + \int_{|z| > \frac{1}{2}d(x)} \hspace{-5pt} \frac{v(x) - v(x-z)}{|z|^{N+2s}}\ dz\nonumber\\
        &=: I_1 + I_2 \text{ (say)}.
    \end{align}
    For $I_1$, Taylor's expansion yields $\theta_1, \theta_2 \in (0,1)$ such that
    \begin{align*}
        |I_1| &= \left|\frac{1}{4} \int_{|z| \leq \frac{1}{2}d(x)} \frac{z^T D^2v (x+ \theta_1 z) z + z^T D^2v (x- \theta_2 z) z}{|z|^{N+2s}}\ dz \right|\\
        &\leq \frac{k_3}{4} \int_{|z| \leq \frac{1}{2}d(x)} \left( \frac{1}{d(x+ \theta_1 z)} + \frac{1}{d(x- \theta_2 z)} \right) |z|^{-N-2s+2}\ dz, \text{ by \eqref{Hessian_estimate}}.
    \end{align*}
    Note that for every $x\in \Om_{\eta}$ and $|z| \leq \frac{1}{2}d(x)$, the points $x\pm \theta_i z \in \Om$ and $d(x\pm \theta_i z) \geq \frac{1}{2}d(x)$, for $i=1,2$. As a result, we obtain the following bound for $I_1$:
    \begin{equation}\label{est_1}
        |I_1| \leq \frac{k_3}{d(x)} \int_0^{\frac{1}{2}d(x)} r^{-2s+1}\ dr \leq \frac{k_4}{d(x)^{2s-1}}, \text{ for } s\in \left[\frac{1}{2},1\right).
    \end{equation}

    To estimate $I_2$, we split the argument into two cases. When $s\in \left(\frac{1}{2},1\right)$, a direct calculation gives
    \begin{equation}\label{est_2}
        |I_2| \leq \|v\|_{C^1(\sr^N)} \int_{|z|>\frac{1}{2}d(x)} |z|^{-N-2s+1}\ dz \leq \frac{k_5}{d(x)^{2s-1}}.
    \end{equation}
    In the case $s= \frac{1}{2}$, since $\Om$ is a bounded domain and thanks to the construction of $\tilde{e}_0$ as in lemma \ref{extension of e}, we can choose $R>0$ large, independent of $x\in\Om_{\eta}$, such that whenever $|z|\geq R$ we have $v(x-z)=0$. Hence, in this case $I_2$ can be estimated as follows:
    \begin{align}\label{est_3}
        |I_2| &\leq \int_{\frac{1}{2}d(x) < |z| < R} \frac{|v(x) - v(x-z)|}{|z|^{N+1}}\ dz + \int_{|z|\geq R} \frac{|v(x)|}{|z|^{N+1}}\ dz\nonumber\\
        &\leq \|v\|_{C^1 (\sr^N)} \int_{\frac{1}{2}d(x) < |z| <R} |z|^{-N}\ dz + \|v\|_{L^{\infty}} \int_{|z|\geq R} |z|^{-N-1}\ dz \leq k_6.
    \end{align}
    Substituting \eqref{est_1}, \eqref{est_2} and \eqref{est_3} in \eqref{est_nonlocal}, the claim follows.

    Combining this estimate with \eqref{est_k1}, we thus obtain
    \begin{equation}\label{v_equn_second}
        -\Delta v + (-\De)^s v + \frac{k}{d(x)^{\ba +1}} v \leq - \frac{k_1}{d(x)^{2-\gamma}} + \frac{k_2}{d(x)^{2s-1}} \text{ in } \Om_{\eta}, \text{ for } s\in \left[\frac{1}{2},1 \right).
    \end{equation}
    Now choose $\gamma \in \left(1, \text{min} \left\{ 2-\beta, 3-2s \right\} \right)$ if $s\in \left(\frac{1}{2},1 \right)$ and $\gamma \in (1, 2- \ba)$ if $s= \frac{1}{2}$ so that for any $s\in \left[ \frac{1}{2}, 1 \right)$ we have
    \begin{equation}\label{v_equn_third}
        -\Delta v + (-\De)^s v + \frac{k}{d(x)^{\ba +1}} v \leq 0 \text{ in } \Om_{\eta},
    \end{equation}
    provided $\eta>0$ is small enough. The above equation is analogous to equation \eqref{v_equn} in the proof of lemma \ref{hopf_first}. Repeating a similar argument and applying the weak comparison principle, we conclude that $\frac{\pa w}{\pa \nu}<0$, or $\frac{\pa u_2}{\pa\nu} < \frac{\pa u_1}{\pa\nu}$ on $\pa\Om$. This completes the proof in the case $s\in \left[ \frac{1}{2},1 \right)$.
\end{proof}

\end{document}
